\section{Randomization Integrity and Sensitivity Analyses}
\subsection{Allocation and exposure fidelity}
Eligibility and washout status will be logged before assignment. The randomization algorithm will preserve the planned allocation probabilities conditional on the prespecified participant/session strata. Assignment, successful playback, headphone configuration, channel balance, interruption, and child-requested stop will be separate fields. Auditory level and routing will be checked on qualified configurations before the pediatric study.

\subsection{Estimand and missingness sensitivity}
The primary assignment-based A--B comparison is the participant-weighted restricted mean recovery-time contrast through the common horizon defined above. Safety-related stops are treated using the prespecified composite failure strategy; missing telemetry and broken recovery measurement contexts are not automatically classified as failures. Sensitivity bounds will vary the unknown recovery status of missing events by arm and assess how much differential missingness would reverse the substantive conclusion. An as-delivered contrast is descriptive because adherence after randomization can be informative.

\subsection{Carryover and competing explanations}
Sequence, previous condition, elapsed washout, session, and exposure count will be documented before analysis. A prespecified first-exposure comparison will probe carryover without replacing the full randomized analysis. The secondary B--C contrast distinguishes general auditory pacing from the alternating-side hypothesis. Any A--B advantage accompanied by increased discomfort, distraction, or stopping will be interpreted as a trade-off, not an unqualified benefit.

\subsection{Falsification}
A bilateral-specific mechanism is not supported if A and B have comparable proximal outcomes, if the apparent effect depends on an unplanned subset, or if acoustical mismatch prevents attribution to alternation. Precision insufficient to rule out both meaningful benefit and harm is inconclusive. Neither the stimulation procedure nor its outcome is equated with EMDR or trauma treatment.